\section{Related work}
\label{related-work}

\subsection{From quality charts to product graphs}

Structured descriptions that link what customers want to what a product does and what it is made of are older than computing in design. Quality function deployment began in Japanese manufacturing in the late 1960s, was taken up by Toyota and its suppliers in the 1970s, and reached Western industry as the ``house of quality'' \citep{ed2024,Hauser1988,Chan2002}. Its matrices relate customer requirements to engineering characteristics, and then to parts and processes, so that every design decision can be traced to a need. Related schemes followed, each making a different link explicit: morphological analysis lays out the alternatives for each function \citep{Zwicky1967,Ritchey2006,Ritchey2018}; the design structure matrix records the dependencies between components \citep{Steward1981,Browning2001,Eppinger2012}; function--behaviour--structure separates what a design is for from what it is \citep{Gero1990,Gero2004}; the functional basis standardises the vocabulary of functions and flows \citep{Stone2000,Hirtz2002}; and technology roadmaps link market needs to products and technologies over time \citep{Phaal2004}. The innovation matrix used here \citep{Tchirsky,berengueres2014} combines several of these links in one layered graph, from market needs through functions and components to technologies and scientific fields (it is distinct from the design-led innovation matrix of \citet{Kyffin2009}). All of these were made for people, typically a team in front of a wall chart, who read them selectively.

\subsection{LLMs that innovate with design methods}



Language models were first used as concept generators \citep{Zhu2022,Ma2023} and compared with people: their concepts are less diverse than crowdsourced ones \citep{Ma2025}, their research ideas are judged more novel but slightly less feasible than experts' \citep{Si2024}, and working with them raises individual creativity while narrowing collective diversity \citep{Doshi2024,Anderson2024,Ashkinaze2025,Wadinambiarachchi2024}. Large language models are often conceptualized as a ``lossy compression of the web'' \citep{Chiang2023}, distilling vast online corpora into compact neural weights \citep{Karpathy2023}. Recent systems give the model the structured methods above. An FBS ontology steers concept generation \citep{Chen2024a}; generated text and images enrich concepts \citep{Chen2024}; 

TRIZ is applied by single models and by agent teams \citep{Chen2024c,Jiang2025,Szczepanik2025,Guo2026} and compared with human TRIZ practice on industrial problems \citep{Livotov2026}, and problems are represented as graphs for the model to work on \citep{Witkowicz2026,Sun2026};
 morphological analysis is automated \citep{Chen2024b}; and agentic workflows chain whole design processes \citep{Ding2023,Massoudi2026}. iDesignGPT \citep{Liu2026} goes furthest: agent clusters run problem definition with QFD and Kano analysis, then functional decomposition, morphological analysis, SCAMPER \citep{Eberle1971} and TRIZ, and produce a product design specification.

\subsection{What ``innovative'' is checked against}

These systems claim to innovate, but the claim usually rests on ratings: experts or a model score novelty and originality, or embedding distances measure how far the generated concepts spread from each other. iDesignGPT, for example, reports that it ``surpasses benchmark models in generating innovative, modular, and rational solutions'', on the basis of expert scores against GPT-4o baselines and semantic novelty metrics; patents are searched as an input to the workflow, but the generated concepts are not compared with prior art \citep{Liu2026}. This is typical. Novelty is judged relative to other generated ideas or to a rater's memory, seldom against what already exists, although patent-based measures of technological novelty are well established \citep{Fleming2001,Verhoeven2016,Uzzi2013} and patent-derived knowledge networks are available for design \citep{Sarica2020,Luo2021}. An idea a rater finds novel may be on the market already. Our own study shares this limitation: its judge rates novelty from knowledge, without a prior-art search, and several of its top-rated ideas, such as solid-state RF heating and RF thawing, already exist in industry, as the judge itself noted. A second gap is cost: none of these systems reports what its outputs cost, so a more elaborate pipeline cannot be weighed against a simpler one.

\subsection{Context and cost in LLMs}

More context is not free even when it fits in the maximum context window available. In practice, models exhibit an ``effective context window'' far smaller than their advertised capacity, with performance degrading sharply as input token length increases \citep{Paulsen2026}, a limitation driven by positional frequency biases and attention sparsity that restrict effective context scaling \citep{An2024}. Models use information in the middle of long inputs poorly \citep{Liu2024}, their reasoning degrades as inputs grow long well before the limit \citep{Levy2024}, and they struggle when relevant content must be inferred rather than matched \citep{Modarressi2025}; retrieval can reach results comparable to long context at a fraction of the cost \citep{Li2024}. Elsewhere, cost is becoming a first-class outcome: cascades and routers match strong models at lower cost \citep{Chen2024d,Ong2024,Ding2024}, reasoning strategies are compared at equal token budgets \citep{Wang2024,Han2025}, and agent benchmarks report accuracy together with cost \citep{Kapoor2025,Wang2025}. This work concerns tasks with right answers. Ideation has none, and its efficiency has rarely been measured.

\subsection{Measuring ideation}

Ideation is usually assessed by quantity, quality, novelty and variety \citep{Shah2003}, or by novelty and usefulness together \citep{Sarkar2011}. But a brainstorm is worth its best ideas, since the rest are discarded: idea-generation processes should be compared on the quality of the best idea, not the average \citep{Girotra2010}, and extreme-value models describe how the best of many concepts improves with the number tested \citep{Dahan2001}. Mean scores, the usual summary in LLM ideation studies, hide exactly this difference. We combine the tail and the cost in one measure, top ideas per token, and vary only how the product is presented, holding the model, the method and the prompt fixed.
